% Chapter 2; source PDF: references/chapman-rizzoli-1989/ChapmanRizzoli0_2.pdf.
% Source-page comments preserve printed/physical page correspondence.

\setcounter{chapter}{1}
\sourcesetpage{18}

% -----------------------------------------------------------------------------
% Source printed page 18 / physical page 28
% -----------------------------------------------------------------------------
\chapter{Acoustic waves}

Being now equipped with some ideas about wave motions, it is useful to consider an
example of waves which occurs in both the ocean and the atmosphere and which can
illustrate many of the ideas in a rather simple way. Acoustic or sound waves, as
Lighthill (1978) points out, are the most fundamental waves in fluids because they can
exist in the absence of any external force field. Instead of gravity or rotation, for
example, providing a restoring force for the motions, the restoring force for acoustic
waves is the fluid's resistance to compression (i.e., its compressibility).

\section{Basic physics}

When viscous dissipation, rotation and gravitational forces are neglected, the
momentum and continuity equations are
\begin{align*}
	\frac{\partial\vec u^{\,*}}{\partial t}
	+\vec u^{\,*}\cdot\nabla\vec u^{\,*}
	                                 & =-\frac{1}{\rho^*}\nabla p^*, \\
	\frac{\partial\rho^*}{\partial t}
	+\nabla\cdot(\rho^*\vec u^{\,*}) & =0.
\end{align*}

\sourcepagebreak

% -----------------------------------------------------------------------------
% Source printed page 19 / physical page 29
% -----------------------------------------------------------------------------
An equation relating density and pressure may be obtained from the first law of
thermodynamics (Batchelor, 1967; Chapter 3). It can be shown that, if
$\rho^*=\rho^*(p^*,T)$ and the motions are adiabatic so that $\partial S/\partial t=0$ where $S$ is
the entropy, then $\rho^*=\rho^*(p^*,S)$ and
\[
	\left(\frac{D\rho^*}{Dt}\right)_S
	=\left(\frac{\partial\rho^*}{\partial p^*}\right)_S
	\left(\frac{Dp^*}{Dt}\right)_S.
\]
A solution of these equations, although trivial, is
\[
	\rho^*=\rho_0;\qquad p^*=p_0;\qquad \vec u^{\,*}=0.
\]
This solution is not very exciting, so we would like to study small deviations from it.
Thus, we write
\[
	\rho^*=\rho_0+\rho;\qquad
	p^*=p_0+p;\qquad
	\vec u^{\,*}=0+\vec u
\]
where $\rho,p,\vec u$ are of infinitesimal amplitude. After substituting into the original
equations and neglecting products of small quantities, we have
\begin{align*}
	\frac{\partial\vec u}{\partial t}                       & =-\frac{1}{\rho_0}\nabla p,           \\
	\frac{\partial\rho}{\partial t}+\rho_0\nabla\cdot\vec u & =0,                                   \\
	\frac{\partial\rho}{\partial t}                         & =c^{-2}\frac{\partial p}{\partial t},
\end{align*}
where
\[
	c^2=\left(\frac{\partial\rho^*}{\partial p^*}\right)_S^{-1}.
\]
After eliminating $\vec u$ and $\rho$ in favor of $p$, we obtain
\begin{waveequation}
	\frac{\partial^2p}{\partial t^2}-c^2\nabla^2p=0.
\end{waveequation}
We recognize this as a wave equation which was listed in Chapter 1. It is easy to show
that the other variables $\rho,u,v,w$ each satisfy a similar equation.

\sourcepagebreak

% -----------------------------------------------------------------------------
% Source printed page 20 / physical page 30
% -----------------------------------------------------------------------------
\section{Plane waves}

Consider a homogeneous medium; $c(\vec x,t)=c_0$. Then
$p=a e^{-i\sigma t+ikx+i\ell y+imz}$ solves the wave equation provided
\begin{waveequation}
	\sigma^2=c_0^2(k^2+\ell^2+m^2)
\end{waveequation}
which is the dispersion relation. For fixed $\sigma$, the locus of allowed wavenumbers in
$k,\ell,m$ space is a sphere of radius $\sigma/c_0$. All wavenumbers
$\vec k=k\hat i+\ell\hat j+m\hat k$ extending from the center of this sphere to its
surface are allowed. In a given plane wave, phases propagate along the wavenumber
vector at speed $c_0$; that is, $\sigma/|\vec k|=c_0$, so the waves are nondispersive.
The group velocity $\vec c_g$ is defined by
\[
	c_{gx}=\frac{\partial\sigma}{\partial k};\qquad
	c_{gy}=\frac{\partial\sigma}{\partial\ell};\qquad
	c_{gz}=\frac{\partial\sigma}{\partial m}
\]
and it is easy to show that $|\vec c_g|=c_0$.

If $p=a e^{i(\vec k\cdot\vec x-\sigma t)}$, then the momentum equations say
$-i\sigma\vec u=-i\vec k p/\rho_0$, or
\[
	\vec u=\frac{\vec k a}{\sigma\rho_0}
	e^{i(\vec k\cdot\vec x-\sigma t)}.
\]
This means that $\vec u$ and $\vec k$ are parallel, i.e. these are longitudinal waves
(displacement is parallel to the direction of wave propagation). Also, $\vec u$ and $p$
are in phase in this travelling plane wave.

\section{Reflection at a solid boundary}

Suppose a plane wave of the form
\[
	p_{inc}=p_0e^{-i\sigma t+ikx+i\ell y-imz}
\]

\sourcepagebreak

% -----------------------------------------------------------------------------
% Source printed page 21 / physical page 31
% -----------------------------------------------------------------------------
is incident upon a solid boundary. At the solid boundary, the normal velocity must
vanish; $\vec u\cdot\hat n=0$ which means that $\nabla p\cdot\hat n=0$. If the solid
boundary is at $z=0$, then the boundary condition is
\[
	p_z=0\qquad\text{at }z=0.
\]

% Vector reconstruction; source PDF ChapmanRizzoli0_2.pdf, physical page 31, trim=105bp 444bp 100bp 180bp.
\wavevectorart[0.72]{ch02-p021-solid-boundary-reflection}

To satisfy this boundary condition, we must add a reflected wave
\[
	p_{ref}=p_0e^{-i\sigma t+ikx+i\ell y+imz}
\]
to the incident wave. The solution is
\[
	p=p_{inc}+p_{ref}=2p_0e^{-i\sigma t+ikx+i\ell y}\cos mz.
\]

Suppose the solid boundary is tilted, say $z=\alpha x$, and the incident energy
approaches along $\vec k_i$ while the reflected energy travels along $\vec k_r$.

\sourcepagebreak

% -----------------------------------------------------------------------------
% Source printed page 22 / physical page 32
% -----------------------------------------------------------------------------
% Vector reconstruction; source PDF ChapmanRizzoli0_2.pdf, physical page 32, trim=140bp 512bp 135bp 91bp.
\wavevectorart[0.67]{ch02-p022-specular-reflection}

If $p_{inc}$ and $p_{ref}$ are to sum such that $p=p_{inc}+p_{ref}$ satisfies
$\partial p/\partial n=0$ at the solid boundary, then we must have
\[
	|\vec k_i|\cos\theta_i=|\vec k_r|\cos\theta_r,
\]
i.e., the projection of the incident wavenumber on the boundary must equal the
projection of the reflected wavenumber on the boundary. But we know that
$|\vec k_i|=|\vec k_r|=\sigma/c_0$, so $\theta_r=\theta_i$. That is, the reflection of
these waves is specular. (This is not true of all waves, however.) Note that
$\partial p/\partial n=0$ at the boundary means $\vec u\cdot\hat n=0$ there, so
$\vec u_{inc}\cdot\hat n=-\vec u_{ref}\cdot\hat n$.

\section{Plane waves in a channel}

A very important aspect of wave motion is the effect of boundaries which form a
channel or \emph{waveguide}. Thus far, the plane waves we have considered have not been
restricted in the choice of wavenumbers. That is, the entire continuum of $k,\ell,m$
choices has been available, provided we were willing to accept whatever frequency was
required by the dispersion relation. We saw that the form of the plane wave was
altered somewhat due to the presence of one boundary, so now we consider the effect of
a second boundary.

\sourcepagebreak

% -----------------------------------------------------------------------------
% Source printed page 23 / physical page 33
% -----------------------------------------------------------------------------
% Vector reconstruction; source PDF ChapmanRizzoli0_2.pdf, physical page 33, trim=100bp 420bp 90bp 95bp.
\wavevectorart[0.70]{ch02-p023-waveguide-boundary-problem}

Now the field equation is still valid in the interior of the channel, but the free waves
must satisfy $\partial p/\partial z=0$ on both boundaries, at $z=0,-D$. To find a
solution, we assume that the waves are free to travel along the channel, but that the
cross-channel dependence is unknown.
\[
	p(x,y,z,t)=p_0e^{-i\sigma t+ikx+i\ell y}P(z).
\]
This is substituted into the field equation to obtain an equation for the cross-channel
structure
\[
	P_{zz}+(\sigma^2/c_0^2-k^2-\ell^2)P=0
\]
\[
	P_z=0\qquad\text{at }z=0,-D.
\]
This equation has the solution
\[
	P(z)=\cos n\pi z/D
\]
provided that
\begin{waveequation}
	\sigma^2/c_0^2=k^2+\ell^2+n^2\pi^2/D^2,
	\qquad n=0,1,2,\ldots
\end{waveequation}

\sourcepagebreak

% -----------------------------------------------------------------------------
% Source printed page 24 / physical page 34
% -----------------------------------------------------------------------------
Notice that these solutions are each a sum of two plane waves
\[
	\frac12p_0e^{-i\sigma t+ikx+i\ell y+in\pi z/D}
	+\frac12p_0e^{-i\sigma t+ikx+i\ell y-in\pi z/D}
\]
which satisfy $\partial p/\partial z=0$ at $z=0$ regardless of whether $n$ is an integer
or not. However, to satisfy $\partial p/\partial z=0$ at $z=-D$, we need
$n=0,1,2,\ldots$. These solutions are called \emph{waveguide modes}.

Another important point to notice is that each three-dimensional plane wave by itself
satisfies the dispersion relation, so that both waves are the usual nondispersive plane
waves if we think of $n$ as continuously variable. Yet the solution viewed as a
two-dimensional plane wave restricted to the channel direction is dispersive!
\[
	c_{ph}=\sigma/(k^2+\ell^2)^{1/2}
	=\pm c_0\left[1+\frac{n^2\pi^2}{(k^2+\ell^2)D^2}\right]^{1/2}.
\]
The horizontal group velocity $\partial\sigma/\partial k$, $\partial\sigma/\partial\ell$ is
\[
	\vec c_g
	=c_0\frac{k\hat i+\ell\hat j}
	{(k^2+\ell^2+n^2\pi^2/D^2)^{1/2}}.
\]
It is parallel to the horizontal wavenumber $k\hat i+\ell\hat j$ but not equal to the
phase velocity.

% Vector reconstruction; source PDF ChapmanRizzoli0_2.pdf, physical page 34, trim=145bp 125bp 170bp 455bp.
\wavevectorart[0.58]{ch02-p024-waveguide-dispersion}

The $n=0$ mode actually is nondispersive.

If we fix the horizontal wavelength $2\pi/(k^2+\ell^2)^{1/2}$, say by a wavemaker of
fixed size perhaps but variable frequency, then there is an infinity of waveguide modes

\sourcepagebreak

% -----------------------------------------------------------------------------
% Source printed page 25 / physical page 35
% -----------------------------------------------------------------------------
$n=0,1,2,\ldots$ of ever increasing frequency
$\sigma^2=c_0^2(k^2+\ell^2+n^2\pi^2/D^2)$. However, if we fix the frequency, then
\[
	(k^2+\ell^2)=\sigma^2/c_0^2-n^2\pi^2/D^2
\]
and only for $n=0,1,\ldots n_{max}$ will $k^2+\ell^2>0$ where
$n_{max}=\operatorname{int}[(D/\pi)(\sigma/c_0)]$. That is, only for
$n=0,1,\ldots n_{max}$ will the waveguide modes propagate down the channel! For
example, consider waves in the $x$-direction only
\begin{align*}
	n<n_{max}\qquad
	p & =e^{-i\sigma t+i(\sigma^2/c_0^2-n^2\pi^2/D^2)^{1/2}x}
	\cos\frac{n\pi z}{D},                                     \\
	n>n_{max}\qquad
	p & =e^{-i\sigma t-(n^2\pi^2/D^2-\sigma^2/c_0^2)^{1/2}x}
	\cos\frac{n\pi z}{D}.
\end{align*}
The first set represents travelling waves. The second set represents \emph{evanescent} waves
which decay exponentially away from their source. Practically, this means that if we
have a harmonic wavemaker in the channel, then we may expect to see more
cross-channel structure near the wavemaker than far away from it.

\section{Scattering at a discontinuity}

We have considered the effect of a solid boundary on the propagation of sound waves.
Suppose, however, that a plane wave encounters a boundary between two fluids at which
the properties change abruptly, i.e. a discontinuity.

\sourcesetpage{26}

% -----------------------------------------------------------------------------
% Source printed page 26 / physical page 36
% -----------------------------------------------------------------------------
% Vector reconstruction; source PDF ChapmanRizzoli0_2.pdf, physical page 36, trim=80bp 515bp 75bp 55bp.
\wavevectorart[0.84]{ch02-p026-interface-scattering}

This discontinuity could represent the air-sea interface or the ocean bottom (which is
not truly a solid boundary because it transmits sound waves). In both cases the
incident wave approaches the discontinuity while travelling through the medium which
has density $\rho_1$ and phase speed $c_1$. The density of the medium on the other side is $\rho_2$
while the phase speed is $c_2$.

For the case on the left (upward propagating incident wave), the incident,
reflected and transmitted waves have the following forms;
\begin{align*}
	p_I & =a e^{-i\sigma t+ikx+im_1z},  \\
	p_R & =Ra e^{-i\sigma t+ikx-im_1z}, \\
	p_T & =Ta e^{-i\sigma t+ikx+im_2z},
\end{align*}
where $R$ is the reflection coefficient and $T$ is the transmission coefficient. Notice that
the incident and reflected waves have the same wavenumber component in $z$ but that
they propagate in opposite directions. The transmitted wave has a different
wavenumber in $z$ because the medium has different properties. The wavenumber in the
direction of the boundary $x$ as well as the frequency $\sigma$ are the same for all three waves
because there is nothing in the fluids which would change them.

\sourcepagebreak

% -----------------------------------------------------------------------------
% Source printed page 27 / physical page 37
% -----------------------------------------------------------------------------
To solve the problem, we require that the pressure as well as the velocity normal
to the boundary $w$ be continuous across the boundary. That is
\begin{align*}
	p_I+p_R                         & =p_T                    &  & \text{at }z=0, \\
	\frac{1}{\rho_1}(p_{Iz}+p_{Rz}) & =\frac{1}{\rho_2}p_{Tz} &  & \text{at }z=0.
\end{align*}
Now, substituting the expressions for $p_I,p_R$ and $p_T$, we obtain
\begin{align*}
	1+R                     & =T,                   \\
	\frac{m_1}{\rho_1}(1-R) & =\frac{m_2}{\rho_2}T.
\end{align*}
From the dispersion relation, $m=\sigma\cos\theta/c$ which changes the second matching
condition to
\[
	\frac{1}{\rho_1c_1}(1-R)\cos\theta_I
	=\frac{1}{\rho_2c_2}T\cos\theta_T.
\]
These can be combined to yield
\begin{waveequation}
	R=\frac{\rho_2c_2\cos\theta_I-\rho_1c_1\cos\theta_T}
	{\rho_2c_2\cos\theta_I+\rho_1c_1\cos\theta_T}
\end{waveequation}
\begin{waveequation}
	T=\frac{2\rho_2c_2\cos\theta_I}
	{\rho_2c_2\cos\theta_I+\rho_1c_1\cos\theta_T}.
\end{waveequation}
Identical expressions for $R$ and $T$ result for the downward propagating incident wave.

We see from these expressions that if the density times the phase speed of the
second medium is much less than that of the first, $\rho_2c_2\ll\rho_1c_1$, then the transmission
coefficient vanishes and the reflection coefficient goes to unity, $T\to0,R\to-1$. This is
consistent with the result we obtained for a solid boundary. It is also nearly the case
for the boundary between the ocean and the atmosphere where $\rho c$ is about $1.5\times10^6$ kg
m$^{-2}$ s$^{-1}$ for the ocean and 400 kg m$^{-2}$ s$^{-1}$ for the atmosphere. So, very little sound is
transmitted from the ocean to the atmosphere. On the other hand, a sound wave in
the atmosphere is actually amplified upon encountering the ocean. That is, if medium

\sourcepagebreak

% -----------------------------------------------------------------------------
% Source printed page 28 / physical page 38
% -----------------------------------------------------------------------------
1 is the atmosphere, then $T\to2$. Of course, the sound wave in the atmosphere travels
so slowly relative to the ocean that its energy flux is generally fairly small, so the
amplification is a rather small effect as well. In either case, the energy flux in the $z$
direction is conserved because
\[
	|p_Iw_I|=|p_Rw_R|+|p_Tw_T|.
\]

To complete the calculation, we must find the angle of the transmitted wave, $\theta_T$.
This is found by writing the frequency on both sides of the discontinuity as
\[
	\sigma=c_1(k^2+m_1^2)^{1/2}=c_2(k^2+m_2^2)^{1/2}.
\]
We can write this in terms of the wave angles since $(k^2+m^2)^{1/2}=k/\sin\theta$. Thus,
\begin{waveequation}
	\frac{\sin\theta_I}{c_1}=\frac{\sin\theta_T}{c_2}
\end{waveequation}
which is known as \emph{Snell's Law}. From this we see that, if $c_1<c_2$, then there exists a
critical angle of incidence
\begin{waveequation}
	\theta_{Ic}=\sin^{-1}(c_1/c_2)
\end{waveequation}
beyond which there is total reflection of the incident wave despite the fact that the
second medium can support sound waves. The boundary is then effectively solid.

% Vector reconstruction; source PDF ChapmanRizzoli0_2.pdf, physical page 38, trim=80bp 95bp 75bp 535bp.
\wavevectorart[0.82]{ch02-p028-total-internal-reflection}

This is called total internal reflection.

\sourcepagebreak

% -----------------------------------------------------------------------------
% Source printed page 29 / physical page 39
% -----------------------------------------------------------------------------
\section{Generation of plane waves}

At this point, it is natural to ask how these plane waves may be generated. What
initial or boundary conditions or forcing terms are needed to generate solutions of the
wave equation corresponding to some physical situation? If there are wavemakers in
the medium, they can be modelled by body forces $\vec F/\rho_0$ and mass sources $Q$
\begin{align*}
	\vec u_t                       & =-\nabla p/\rho_0+\vec F/\rho_0, \\
	\rho_t+\rho_0\nabla\cdot\vec u & =Q.
\end{align*}
Combining these with $p_t=c^2\rho_t$ yields
\begin{waveequation}
	p_{tt}-c^2\nabla^2p=c^2(Q_t-\nabla\cdot\vec F).
\end{waveequation}
We can now consider two types of problems: initial value problems and those forced
from rest. In both types we solve the homogeneous wave equation while satisfying
$\partial p/\partial n=0$ on the solid boundaries and requiring outgoing waves at infinity, i.e. a
\emph{radiation condition}. For the initial value problems, $p$ and $p_t$ are specified at time $t=0$,
while for those forced from rest they are set to zero. Of course, there is not really a
fundamental distinction because solutions of one type may be linearly superposed to
obtain solutions to the other type. The solution procedures may, however, be quite
different.

\subsection{An initial value problem}

Let us consider a one-dimensional initial value problem
\[
	p_{tt}-c^2p_{xx}=0\qquad -\infty<x<\infty
\]
\[
	p(x,0)=P_0(x)\ ;\qquad p_t(x,0)=Q_0(x).
\]

\sourcepagebreak

% -----------------------------------------------------------------------------
% Source printed page 30 / physical page 40
% -----------------------------------------------------------------------------
We will solve this by the \emph{method of characteristics}. The most general solution is
\[
	p=f(x-ct)+g(x+ct).
\]
To satisfy the initial conditions
\begin{align*}
	f(x)+g(x)      & =P_0(x), \\
	-cf'(x)+cg'(x) & =Q_0(x).
\end{align*}
The second integrates to
\[
	f(x)-g(x)=-\frac1c\int_0^x Q_0(x')\,dx'+K
\]
whence
\begin{align*}
	2f(x) & =P_0(x)-\frac1c\int_0^x Q_0(x')\,dx'+K, \\
	2g(x) & =P_0(x)+\frac1c\int_0^x Q_0(x')\,dx'-K.
\end{align*}
These give the solution as
\[
	p(x,t)=\frac12\left[P_0(x-ct)+P_0(x+ct)
		+\frac1c\int_{x-ct}^{x+ct}Q_0(x')\,dx'\right].
\]
Note that $p(x,t)$ depends only on the initial conditions over the range $x\pm ct$.

If $Q_0(x)=0$, then the solution is very simple
\[
	p(x,t)=\frac12[P_0(x-ct)+P_0(x+ct)]
\]
for which case the solution could have been obtained using the Fourier method,
although it is not the method of choice in this problem. Set
\[
	p(x,t)=\int_{-\infty}^{\infty}\bar p(k,t)e^{ikx}\,dk.
\]
Then
\[
	\bar p_{tt}+c^2k^2\bar p=0
\]
\[
	\bar p(k,0)=\bar P_0(k)\ ;\qquad \bar p_t(k,0)=0.
\]

\sourcepagebreak

% -----------------------------------------------------------------------------
% Source printed page 31 / physical page 41
% -----------------------------------------------------------------------------
The solution to this problem is
\[
	\bar p(k,t)=\bar P_0(k)\cos(ckt)
\]
from which
\begin{align*}
	p(x,t)
	 & =\int_{-\infty}^{\infty}\bar P_0(k)\cos(ckt)e^{ikx}\,dk \\
	 & =\frac12\int_{-\infty}^{\infty}\bar P_0(k)
	\left(e^{ikx+ickt}+e^{ikx-ickt}\right)\,dk                 \\
	 & =\frac12[P_0(x-ct)+P_0(x+ct)].
\end{align*}
The integration is trivial in this case but not always.

\subsection{Forcing from rest}

Assume that the forcing has the rather simple form
\[
	Q_t-\nabla\cdot\vec F=\delta(x)q_t(t)
\]
where $q(t)=q_t(t)=0$ for $t<0$ and $q_t$ is finite. Now we solve
\[
	p_{tt}-c^2p_{xx}=\delta(x)q_t(t)c^2
\]
\[
	p(x,0)=p_t(x,0)=0\qquad\text{at }t=0.
\]
We may put the forcing into the boundary condition by
\[
	\int_{0^-}^{0^+}(p_{tt}-c^2p_{xx})\,dx
	=-c^2p_x\big|_{0^-}^{0^+}=c^2q_t(t).
\]
That is, $p_x(x=0+,t)-p_x(x=0-,t)=-q_t(t)$ so that the forcing at $x=0$ is
interpretable as a specified discontinuity there. So we must solve

\sourcepagebreak

% -----------------------------------------------------------------------------
% Source printed page 32 / physical page 42
% -----------------------------------------------------------------------------
% Vector reconstruction; source PDF ChapmanRizzoli0_2.pdf, physical page 42, trim=90bp 485bp 80bp 90bp.
\wavevectorart[0.78]{ch02-p032-forced-source-jump}

where $p^L$ and $p^R$ are solutions on the left and right of the discontinuity, respectively.
Most generally, $p^L$ and $p^R$ are functions of $x\pm ct$. We write them along with the
requirement of symmetry
\[
	p^R(x,t)=p^L(-x,t)
\]
\[
	p^R(x,t)=f(x-ct)+g(x+ct)
\]
\[
	p^L(x,t)=f(-x-ct)+g(-x+ct).
\]
Imposing the jump condition at $x=0$ yields
\[
	f'(-ct)+g'(ct)+f'(-ct)+g'(ct)=-q_t(t)
\]
but this does not specify $f$ and $g$. To specify them, we must impose a \emph{radiation}
condition, i.e.,
\[
	p^R(x,t)=f(x-ct)\qquad p^R\text{ is all right going waves}
\]
\[
	p^L(x,t)=f(-x-ct)\qquad p^L\text{ is all left going waves}.
\]
Now we have
\[
	2f'(-ct)=-q_t(t)
\]
\[
	2f(-ct)=q(t)c
\]
\[
	f(\tau)=\frac{c}{2}q(-\tau/c).
\]

\sourcepagebreak

% -----------------------------------------------------------------------------
% Source printed page 33 / physical page 43
% -----------------------------------------------------------------------------
from which
\[
	p^R(x,t)=\frac{c}{2}q(-x/c+t)
\]
\[
	p^L(x,t)=\frac{c}{2}q(x/c+t).
\]
Thus, forcing at the origin is modelled as a jump in $p_x$ and we must assume that all of
the motion is \emph{away} from the source in order to get a unique answer.

If the forcing were harmonic with $q(t)=e^{-i\sigma t}/(-i\sigma)$ then
\[
	p_{tt}-c^2p_{xx}=c^2\delta(x)e^{-i\sigma t}
\]
and the solution would be
\[
	p^R(x,t)=\frac{-c}{2(i\sigma)}e^{-i\sigma(-x/c+t)}
\]
\[
	p^L(x,t)=\frac{-c}{2(i\sigma)}e^{-i\sigma(x/c+t)}.
\]
In other words, plane waves radiating outwards from $x=0$. The radiation condition
that we imposed models a little bit of dissipation in the sense that the solution looks
dissipationless locally, but nothing is reflected from $|x|\to\infty$ because even small
dissipation attenuates any reflected waves over a long distance. We could, in fact, add
a friction term to the momentum equations and solve again to obtain a solution which
would become the present solution for vanishingly small friction.

\section{Slowly varying medium}

We have considered cases in which the speed of sound remains constant in the medium
or changes abruptly at an interface. However, the speed of sound within the ocean
varies in space because the ocean is not a uniform fluid. In fact the sound speed in the

\sourcesetpage{34}

% -----------------------------------------------------------------------------
% Source printed page 34 / physical page 44
% -----------------------------------------------------------------------------
ocean is sensitive to the temperature, salinity and pressure of the ocean and may be
described by the following empirical formula:
\[
	c(s,T,z)=c_0+\alpha_0(T-10)+\beta_0(T-10)^2+\gamma_0(T-18)^2
	+\delta_0(s-35)+\epsilon_0(T-18)(s-35)+\zeta_0|z|
\]
where the coefficients have the appropriate mks units and have values of
	{\small
		\[
			c_0=1493.0,\quad \alpha_0=3.0,\quad \beta_0=-0.006,\quad \gamma_0=-0.04,
			\quad \delta_0=1.2,\quad \epsilon_0=-0.01,\quad \zeta_0=0.0164.
		\]
	}
This says that the speed of sound varies quadratically with temperature, and linearly
with salinity and depth. The depth effect is due to changes in the ambient pressure.
For typical ocean conditions, the temperature effect dominates in the shallow water,
while the pressure effect dominates in the deep water. The sound speed increases with
an increase in either temperature or depth, so there is typically a sound speed
minimum in the ocean interior.

% Vector reconstruction; source PDF ChapmanRizzoli0_2.pdf, physical page 44, trim=170bp 195bp 170bp 397bp.
% wave-source-mask: pdf=ChapmanRizzoli0_2.pdf; page=44; rect=170bp 382bp 240bp 397bp; origin=lower-left; reason=remove the preceding prose fragment while preserving the complete c label and upper axis
\wavevectorart[0.54]{ch02-p034-sound-speed-profile}

The situation is different in the arctic where there is little effect of warming near the
surface. There, the sound speed tends to decrease right up to the surface.

\sourcepagebreak

% -----------------------------------------------------------------------------
% Source printed page 35 / physical page 45
% -----------------------------------------------------------------------------
% Vector reconstruction of the isolated sound-speed profile artwork.
% Vector reconstruction; source PDF ChapmanRizzoli0_2.pdf, physical page 45, trim=150bp 480bp 140bp 55bp.
\wavevectorart[0.58]{ch02-p035-sound-speed-profiles}

We can examine the effects of these variations in sound speed by applying our
knowledge of ray theory. We must assume that the wavelengths of the acoustic waves
are much less than the scale over which the sound speed changes. That is, the
wavelength must be small compared to the total ocean depth. We will consider only
two dimensions, the vertical and one horizontal. Recalling our discussion of ray theory,
we write the dispersion relation as
\[
	\sigma=\Omega(k,m;z).
\]
Since the medium varies only in $z$, we have $\partial\Omega/\partial x=0$,
$\partial\Omega/\partial t=0$ but $\partial\Omega/\partial z\ne0$.
Thus, the ray equations become
\begin{align*}
	\frac{dN}{dt} & =0,                                  \\
	\frac{dk}{dt} & =0,                                  \\
	\frac{dm}{dt} & =-\frac{\partial\Omega}{\partial z}.
\end{align*}
These say that the component of the wavenumber in the $x$ direction remains constant
in time (which makes sense since the medium varies only in $z$), and that the frequency
remains fixed at the initial frequency. We could integrate these following along a ray

\sourcepagebreak

% -----------------------------------------------------------------------------
% Source printed page 36 / physical page 46
% -----------------------------------------------------------------------------
with the group velocity, but we will examine the qualitative behavior by considering
Snell's Law which can be derived in the same manner as for the case of scattering at
the discontinuity. The frequency can be written in terms of the angle that the
wavenumber makes with the vertical to obtain
\begin{waveequation}
	\frac{\sin\theta_0}{c_0}=\frac{\sin\theta}{c(z)}
\end{waveequation}
or
\[
	\sin\theta=\frac{c(z)}{c_0}\sin\theta_0
\]
where $c_0$ and $\theta_0$ are the initial values.

Consider the case in which the sound speed decreases with depth,
$c=c_0(1+\alpha z)$ (remember that $z$ is positive upwards). This means that a wave
moving upward moves into a region of increasing sound speed, so the angle with the
vertical must increase as well. Thus, the wave moves toward a horizontal path. This
may also be seen from the ray equations where $-\partial\Omega/\partial z<0$, so that
$m$ must decrease with upward motion. Decreasing $m$ leads to a more horizontal
propagation path.

% Vector reconstruction; source PDF ChapmanRizzoli0_2.pdf, physical page 46, trim=80bp 250bp 70bp 472bp.
\wavevectorart[0.82]{ch02-p036-sound-ray-turning}

Similarly if $c$ increases in the deep ocean, sound waves moving downward will be
turned toward the horizontal.

When the ray becomes nearly horizontal, ray theory must be applied very
carefully. From the ray definition, we can write
\[
	\frac{dz}{dx}=\frac{c_{gz}}{c_{gx}}
	=\frac{\partial\Omega/\partial m}{\partial\Omega/\partial k}
	=\frac{m}{k}
	=\left(\frac{\sigma^2}{c^2(z)k^2}-1\right)^{1/2}
\]

\sourcepagebreak

% -----------------------------------------------------------------------------
% Source printed page 37 / physical page 47
% -----------------------------------------------------------------------------
which gives the slope of the ray path. This may be approximated near the critical level
$z_c$ by expanding in a Taylor series to obtain
\[
	\frac{dz}{dx}\simeq
	\left[
		(z-z_c)\frac{d}{dz}\left(\frac{\sigma^2}{c^2(z)k^2}\right)_{z=z_c}
		\right]^{1/2}
\]
which integrates to
\[
	z=z_c+\frac14
	\left[\frac{d}{dz}\left(\frac{\sigma^2}{c^2(z)k^2}\right)\right]_{z=z_c}
	(x-x_0)^2.
\]
Thus, the ray path is parabolic near the critical level, so an upward propagating ray
turns downward.

Similarly, a downward propagating ray which encounters an increasing $c$ at
depth will eventually turn upward (provided it does not intersect the bottom). The
end result is that the minimum in the sound speed acts as a sound channel where
acoustic energy can propagate over hundreds of kilometers without encountering the
bottom provided the incidence angle is not too oblique. Numerous examples are
reproduced in Apel (1987). This is the basis of acoustic tomography, in which this
efficient propagation is used to infer properties of the ocean. Sound waves are
generated at a source and received at a listening station. For a fixed vertical profile of
the sound speed, the rays may be calculated using ray theory. The received signal is
then compared with that expected for a horizontally uniform medium, and differences
are used to deduce various physical phenomena which might have occurred along the
ray paths. This is generally called an \emph{inverse problem} because boundary observations
are used to determine the interior physics, rather than the reverse.
